\documentclass[11pt,a4paper]{article}
\usepackage[T1]{fontenc}
\usepackage[utf8]{inputenc}
\usepackage[main=italian,provide=*]{babel}
\usepackage{amsmath,amssymb}
\usepackage{geometry}
\geometry{margin=2.3cm,top=2.5cm,bottom=2.7cm}
\usepackage{booktabs}
\usepackage[table]{xcolor}
\usepackage{tcolorbox}
\tcbuselibrary{skins,breakable}
\usepackage{titlesec}
\usepackage{enumitem}
\usepackage{needspace}
\usepackage{fancyhdr}
\usepackage{hyperref}
\hypersetup{colorlinks=true,linkcolor=Sepia,citecolor=Sepia,urlcolor=Sepia}
\newcommand{\Tr}{\operatorname{Tr}}
\newcommand{\ket}[1]{\lvert #1 \rangle}
\newcommand{\bra}[1]{\langle #1 \rvert}

% --- palette ---
\definecolor{inkblue}{HTML}{1B3A5C}
\definecolor{lightblue}{HTML}{EAF1F8}
\definecolor{accent}{HTML}{B0562A}
\definecolor{softgray}{HTML}{6B6B6B}
\definecolor{okgreen}{HTML}{2E6B3E}
\definecolor{okgreenbg}{HTML}{EAF5EC}

% --- titoli di sezione ---
\titleformat{\section}
  {\normalfont\Large\bfseries\color{inkblue}}
  {\thesection}{0.6em}{}[{\color{inkblue!40}\titlerule}]
\titlespacing{\section}{0pt}{0.7em}{0.4em}
\titleformat{\subsection}
  {\normalfont\large\bfseries\color{accent}}
  {}{0em}{}
\titlespacing{\subsection}{0pt}{0.5em}{0.2em}

\pagestyle{fancy}
\fancyhf{}
\renewcommand{\headrulewidth}{0.4pt}
\fancyhead[L]{\small\color{softgray}Correlazioni dinamiche sotto rumore --- Passo 4, Parte 2}
\fancyhead[R]{\small\color{softgray}\thepage}
\fancyfoot[C]{\small\color{softgray}Tesi triennale, Samuele --- Relatore: Prof.\ Alessandro Chiesa}

% --- box riutilizzabili ---
\newtcolorbox{keyeq}[1][]{
  colback=lightblue, colframe=inkblue!70, boxrule=0.6pt,
  arc=2pt, left=8pt, right=8pt, top=3pt, bottom=3pt, #1
}
\newtcolorbox{okbox}[1]{
  colback=okgreenbg, colframe=okgreen!60, boxrule=0.6pt,
  arc=2pt, left=7pt, right=7pt, top=2pt, bottom=2pt,
  fonttitle=\bfseries\color{okgreen}, title={#1}
}

\title{\vspace{-1.5em}\color{inkblue}Correlazioni dinamiche sotto rumore --- Passo 4\\[0.2em]
\Large\normalfont\color{softgray}Circuito Hadamard test completo, con errore di lettura}
\author{Note di lavoro --- Tesi triennale, Samuele\\ Relatore: Prof.\ Alessandro Chiesa}
\date{}

\begin{document}
\sloppy
\maketitle
\thispagestyle{fancy}
\vspace{-2.2em}

File: \texttt{correlatori\_rumorosi\_dimero.py}, \texttt{validate\_correlatori\_rumorosi\_dimero.py}.
Punto di lavoro: $b/J=0.35$, $D/J=0.80$, $J=1$, $t=2$.

\section*{Metodologia}
\addcontentsline{toc}{section}{Metodologia}

\begin{itemize}[leftmargin=1.3em,itemsep=0.3em]
\item \textbf{Quattro blocchi, ciascuno transpilato una sola volta}: a
differenza dei Passi 2--3, il circuito Hadamard test ha altri due
sottocircuiti a 2 qubit oltre al passo di Trotter --- il controlled-$W$
(prima dell'evoluzione) e l'anti-controlled-$V$ con la rotazione di base
(dopo). Ciascuno è transpilato singolarmente in base
$\{R_z,\sqrt X,X,\text{CX}\}$ e poi composto, mai transpilato insieme al
resto (stessa cautela già motivata per il Trotter nel Passo 3).
\item \textbf{Errore di lettura, applicato per la prima volta}: i Passi
2--3 non misuravano nulla. Per un readout simmetrico di parametro $p$, la
distribuzione di probabilità letta è $p_0'=(1-p)p_0+p\,p_1$,
$p_1'=(1-p)p_1+p\,p_0$, da cui
\begin{keyeq}
\centering
$\langle Z\rangle_\text{readout} = p_0'-p_1' = (1-2p)\,(p_0-p_1) =
(1-2p)\,\langle Z\rangle_\text{ideale}$
\end{keyeq}
--- una singola moltiplicazione applicata al $\langle Z\rangle$ già
calcolato con il rumore di gate. Scelta analitica, non a shot finiti,
coerente con il resto del progetto (mai rumore di shot altrove).
\end{itemize}

\section*{Verifica}
\addcontentsline{toc}{section}{Verifica}

\begin{okbox}{Tre controlli superati}
\begin{itemize}[leftmargin=1.2em,itemsep=0.2em]
\item \textbf{Limite di rumore nullo} su tre correlatori diversi: identico
a Parte 1 entro $5\times10^{-14}$ (uguale alla precisione macchina, non
solo "vicino").
\item \textbf{Conteggio gate}: incremento di esattamente 3 CNOT per unità
di $N$ (contributo del passo di Trotter), confermato $N=1,\dots,40$.
\item \textbf{Formula di readout verificata indipendentemente}: simulazione
Monte Carlo con \texttt{ReadoutError} vero e $2\times10^5$ shot vs. formula
analitica --- $\langle Z\rangle=+0.044089$ (analitico) contro $+0.046550$
(Monte Carlo), differenza $2.46\times10^{-3}$, entro l'errore statistico
atteso a $3\sigma$ ($\approx6.7\times10^{-3}$).
\end{itemize}
\end{okbox}

\needspace{16\baselineskip}
\section*{Risultato: $C_{21}^{xx}(t=2)$ in funzione di $N$}
\addcontentsline{toc}{section}{Risultato}

Con il rumore di riferimento (\texttt{ibm\_torino}, $p_\text{readout}=2.3\times10^{-2}$
incluso):

\begin{center}
\renewcommand{\arraystretch}{1.2}
\begin{tabular}{@{}cc@{}}
\toprule
$N$ & $|C_{21}^{xx}(t=2)|$ \\
\midrule
$1$ & $0.102605$ \\
$2$ & $0.119573$ \\
$4$ & $0.238728$ \\
$5$ & $0.250002$ \quad ($\approx N^*$) \\
$6$ & $0.247029$ \\
$10$ & $0.214978$ \\
$20$ & $0.158067$ \\
$40$ & $0.096323$ \\
$80$ & $0.035454$ \\
\bottomrule
\end{tabular}
\end{center}

\begin{okbox}{$N^*$ individuato anche sui correlatori}
Massimo a $N^*=5$ --- più basso dell'$N^*=8$ trovato al Passo 3 sulla
fedeltà. Atteso: il circuito dei correlatori ha più gate a parità di $N$
(due blocchi in più rispetto al solo Trotter), quindi il rumore si
accumula più in fretta e conviene fermarsi prima.
\end{okbox}

\needspace{6\baselineskip}
\section*{Prossimo passo}
\addcontentsline{toc}{section}{Prossimo passo}

Passo 5, l'ultimo: scan sistematico su
$(\varepsilon_{1q},\varepsilon_{2q},p_\text{readout})$ per vedere come si
sposta $N^*$ --- finora misurato solo al punto di riferimento
\texttt{ibm\_torino}, sia sulla fedeltà (Passo 3, $N^*=8$) sia sui
correlatori (qui, $N^*=5$).

\end{document}
